% Section: general real saddle-node unfoldings.  \input from main.tex after Section~\ref{sec:char}.
% Source of the audit: docs/general-statements-zh.md; numerics: docs/generality-check-zh.md,
% docs/quad-kneser-zh.md.
\section{General real saddle-node unfoldings}\label{sec:general}

Sections~\ref{sec:transition}--\ref{sec:char} use the exponential only through a
few constants.  This section states their results for a general real
saddle-node unfolding, lists the changes in the proofs, and adds one exact
invariance statement.  Tetration is the case $a_2=\tfrac12$, $\gamma=1$,
$\rho=\tfrac13$ below.

\subsection{Setting}
Let $f_\mu(w)$ be a family with the following properties.
\begin{enumerate}
\item[(H1)] $f_\mu(w)$ is entire in $w$, holomorphic in $\mu$ near $\mu_0$, and
  real for real $(\mu,w)$.
\item[(H2)] $f_{\mu_0}(w^*)=w^*$, $f_{\mu_0}'(w^*)=1$ and
  $a_2:=f_{\mu_0}''(w^*)/2>0$ (replace $w$ by $-w$ if necessary).
\item[(H3)] $\gamma:=-\partial_\mu f_\mu(w^*)|_{\mu_0}>0$ (reverse $\mu$ if
  necessary).  For $s:=\mu_0-\mu>0$ small there are two real fixed points
  $L<L_2$ near $w^*$, $L$ attracting with multiplier $\lambda$ and $L_2$ repelling
  with multiplier $\lambda_2$.
\item[(H4)] The base point $w_0<w^*$ is real, and $f_{\mu_0}$ is increasing with
  $f_{\mu_0}(w)>w$ on $[w_0,w^*)$.  Then for $s$ small, $w_0<L$ lies in the
  immediate basin of $L$, and the Koenigs map satisfies $\sigma(w_0)<0$.
\item[(H5)] $B_1\ne0$, where $B_n$ are the Fourier coefficients of the inverse upper
  horn map of the germ at $s=0$ (used only where indicated).
\end{enumerate}
Write the germ as $g(u)=f_{\mu_0}(w^*+u)-w^*=u+a_2u^2+a_3u^3+\cdots$, with
iterative residue $\rho=1-a_3/a_2^2$.  Its formal Abel function is
$\alpha(u)=-1/(a_2u)+\rho\log(-u)+\sum_{k\ge1}d_ku^k$, and $\Phi_{\mathrm{att}}$,
$\Phi_{\mathrm{rep}}$, $B_n$ are defined as in Section~\ref{sec:horn}.  The
normalisation $F(0)=1$ becomes $F(0)=w_0$: $\Rr^{-1}(w_0)=0$,
$u^*:=w_0-w^*$ and $a:=\Phi_{\mathrm{att}}(u^*)\in\mathbb R$.  The definitions of
$S$, $T$, $t_n$, $\tau_n$, $h$, $\Lambda$, $\Kk^W$ and $\mathcal C(Y)$ are those of
Sections~\ref{sec:transition} and~\ref{sec:char}, with $\mu$ in place of $b$.

No global hypothesis enters.  All estimates take place in a fixed disc about
$w^*$; (H1) is used only to make $S$ entire.

\begin{proposition}[Scales]\label{prop:gen-scale}
As $s\to0^+$,
\[
 |\log\lambda|,\ \log\lambda_2=2\sqrt{a_2\gamma s}\,\bigl(1+O(\sqrt s)\bigr),\qquad
 \Lambda=\exp\Bigl(-\frac{2\pi^2}{\sqrt{a_2\gamma s}}\bigl(1+o(1)\bigr)\Bigr),
\]
and $p:=|\log\lambda|\log\lambda_2$ is holomorphic in $s$ with
$p=4a_2\gamma s+O(s^2)$.  Moreover, with $A_1=1/\log\lambda$, $A_2=1/\log\lambda_2$
and $u_1=L-w^*$, $u_2=L_2-w^*$,
\[
 A_1+A_2\to\rho,\qquad A_1(u_1-u_2)\to1/a_2,\qquad h_2-h\to2\pi\rho,
\]
where $h_2=2\pi/\log\lambda_2$.
\end{proposition}

\begin{proof}
Write the Weierstrass preparation~\eqref{eq:weier} as
$f(u)-u=(u^2-\sigma_1u+\sigma_2)k_s(u)$ with $k_0(0)=a_2$.  Then
$\sigma_2=f_s(0)/a_2+O(s^2)=-\gamma s/a_2+O(s^2)$, so the discriminant is
$D(s)=4\gamma s/a_2+O(s^2)$, and $\lambda_{1,2}=1\mp a_2\sqrt D+O(D)$.  The
limits of $A_1+A_2$ and $A_1(u_1-u_2)$ are those of Lemma~\ref{lem:res}, with
$1-k'(0)/k(0)^2=1-a_3/a_2^2$ and $1/k(0)=1/a_2$.  Finally
$h_2-h=2\pi(A_1+A_2)$.  For tetration this gives $C_\eta$ of
Proposition~\ref{prop:saddle} and $p=2s+\cdots$.
\end{proof}

\subsection{The theorems}

\begin{theorem}[Confluence, general]\label{thm:gen-A}
Under \textup{(H1)--(H4)}, $\tau_n(\mu)\to B_n\eee^{2\pi\ii na}$ as $\mu\to\mu_0^-$ for every
$n\ge1$, with $\tau_n$ taken on the branch $\operatorname{Im}t_0=h/2$.  Under
\textup{(H5)} also $t_n/t_1^n\to B_n/B_1^n$.
\end{theorem}

\begin{theorem}[Sewing solution, general]\label{thm:gen-B}
Under \textup{(H1)--(H4)} there is $\mu_1<\mu_0$ such that for
$\mu\in(\mu_1,\mu_0)$ the following hold.  The sewn solution $\Kk^W_\mu$,
normalised by $\Kk^W_\mu(0)=w_0$, is the unique element of $\mathcal C_\mu(Y)$.  Its
coefficients satisfy $|c_n/\Lambda^n-\tau_n(\mu)|\le C_n\Lambda$, with $C_n$ as in
Theorem~\ref{thm:W2}(4) and constants $Y_B,Y_0,D,M$ that depend on the family.
The family is holomorphic in $\mu$ near $(\mu_1,\mu_0)$, and for $\operatorname{Im}\mu>0$
small it tends to $L$ and $L_2$ as $\operatorname{Im}z\to+\infty$ and $-\infty$,
respectively.
\end{theorem}

\begin{theorem}[First and higher orders, general]\label{thm:gen-D}
Under \textup{(H1)--(H4)}, and for $n$ with $B_n\ne0$,
\[
 \log\frac{\tau_n}{B_n\eee^{2\pi\ii na}}\sim\sum_{j\ge1}\kappa^{(n)}_jp^j,\qquad
 \kappa^{(n)}_1=\frac{\delta t_n}{4a_2\gamma\,B_n}-\frac{2\pi\ii n\,\delta t_0}{4a_2\gamma},
\]
with $\delta t_m$ defined as in Theorem~\ref{thm:D} from the orbits of the germ $g$
and of its reflected local inverse.  Each $\kappa^{(n)}_j$ is given by convergent
parabolic series as in Theorem~\ref{thm:E}.
\end{theorem}

\begin{proof}[Proof of Theorems~\ref{thm:gen-A}--\ref{thm:gen-D}]
The proofs of Sections~\ref{sec:transition}--\ref{sec:char} apply verbatim,
with the following replacements.
\begin{itemize}
\item In Lemmas~\ref{lem:res}--\ref{lem:petal}, Theorem~\ref{thm:C} and
  Lemma~\ref{lem:finite}, replace
  $\alpha_2=-2/u+\tfrac13\log(-u)$ by $-1/(a_2u)+\rho\log(-u)$ and $-2/u$ by
  $-1/(a_2u)$.  The limiting petal is $\{\operatorname{Re}(-1/(a_2u))>R\}$.
\item In Theorem~\ref{thm:C}(2), the reflected inverse
  $G_\mu(v)=-f_\mu^{-1}(-v)$ is the local inverse near $w^*$.  It is jointly
  holomorphic, its germ is $v+a_2v^2+(2a_2^2-a_3)v^3+\cdots$, and its iterative
  residue is $-\rho$.
\item In Lemma~\ref{lem:realT}, $1<L$ is replaced by \textup{(H4)}: $S$ maps
  $\mathbb R$ onto $(L,L_2)$, where $\sigma>0$, while $\sigma(w_0)<0$.
\item In Lemma~\ref{lem:band}, the bound $\pi/3$ becomes $\pi|\rho|$ and
  $h_2=h+2\pi/3+o(1)$ becomes $h_2=h+2\pi\rho+o(1)$.  Injectivity of $f_\mu$ on
  $D_{r'}$ follows from $f'_{\mu_0}(w^*)=1$ after shrinking $r'$.
\item In Proposition~\ref{prop:weld}, the singular set of $\Rr$ is the closed
  set $\Sigma_\Rr$ of inverse-branch images of the singular values of $f$.
\item In Lemma~\ref{lem:model}, $d_i=\tfrac i2u^{i-1}+\cdots$ becomes
  $d_i=ia_2u^{i-1}+\cdots$, and $D(s)=8s+\cdots$ becomes $4\gamma s/a_2+\cdots$.
  In Theorem~\ref{thm:D}, $p'(0)=2$ becomes $p'(0)=4a_2\gamma$
  (Proposition~\ref{prop:gen-scale}).
\item In Theorem~\ref{thm:W1}(2), $\lambda_{1,2}=1\mp a_2\sqrt{D(s)}+O(D)$ gives
  $\partial_\mu\lambda_1>0>\partial_\mu\lambda_2$, hence
  $\operatorname{Im}\mu>0\Rightarrow\arg\lambda_1>0>\arg\lambda_2$.
\item Statements involving ratios or logarithms use \textup{(H5)} in place of
  Proposition~\ref{prop:horn-nonzero}.  Corollary~\ref{cor:sewn-cusp-local}
  needs, in addition, that $w_0$ is not precritical for $f_{\mu_0}$.
\end{itemize}
Everything else is literally unchanged.  In particular the welding identity,
the Wiener-algebra contraction and the characterisation use only the two
Koenigs charts and the relation $\sigma\circ\Rr(w)=\sigma(w_0)\lambda^w$.
\end{proof}

\begin{remark}[On \textup{(H5)}]\label{rem:gen-B1}
The proof of Proposition~\ref{prop:horn-nonzero} gives $B_1\ne0$ whenever
$f_{\mu_0}$ restricted to its immediate parabolic basin has exactly one singular
value.  This covers $w^2+c$, whose basin contains the single critical point.
For $v-\sin v+c$ the critical values $2\pi k+1$ are real and only $k=0$ lies in
the basin, but the map is not of finite type, so applying the argument requires
a check.  For $w+(w^2-\epsilon)\eee^w$ we only know $B_1\ne0$ numerically.
\end{remark}

\subsection{Independence of the base point}

\begin{proposition}\label{prop:gen-base}
Let $w_0$ and $w_0'$ both satisfy \textup{(H4)}, and mark with a prime the
quantities computed from $w_0'$.  Then for every $\mu<\mu_0$ near $\mu_0$ and every
$n\ge1$,
\[
 \tau_n'(\mu)=\tau_n(\mu)\,\eee^{2\pi\ii n\Delta(\mu)},\qquad
 B_n\eee^{2\pi\ii na'}=B_n\eee^{2\pi\ii na}\,\eee^{2\pi\ii n\Delta_0},
\]
with $\Delta(\mu)=\Rr_\mu^{-1}(w_0')$ and $\Delta_0=\Phi_{\mathrm{att}}(u^{*\prime})-\Phi_{\mathrm{att}}(u^*)$
both real.  Hence $|\tau_n(\mu)|$, $|c_n|$ and $|\hat c_n|$ do not depend on the
base point, and neither does $\operatorname{Re}\kappa^{(n)}_j$ for any $j$.
\end{proposition}

\begin{proof}
By \textup{(H4)}, $\sigma(w_0)$ and $\sigma(w_0')$ are negative, so
$\Delta=\log(\sigma(w_0')/\sigma(w_0))/\log\lambda$ is real.  Then
$\Rr'^{-1}=\Rr^{-1}-\Delta$ and $T'=T-\Delta$, so $t_n'=t_n$ for $n\ne0$ and
$t_0'=t_0-\Delta$.  The branch condition $\operatorname{Im}t_0=h/2$ is preserved,
and $\tau_n'=\tau_n\eee^{2\pi\ii n\Delta}$.  At $\mu_0$, $u^*$ and $u^{*\prime}$ lie
on the real attracting axis, where $\Phi_{\mathrm{att}}$ is real.  The sewn
solutions for the two base points differ by a real translation, which does not
change $|c_n|$ (Remark~\ref{rem:invariance}).  Taking real parts of the logarithm of
the first identity divided by the second proves the last assertion.
\end{proof}

\begin{corollary}\label{cor:gen-invariant}
$|B_n|$ and the real parts $\operatorname{Re}\kappa^{(n)}_j$ are invariant under
orientation-preserving real-analytic conjugacy of the unfolding near
$(\mu_0,w^*)$ and under reparametrisation of $\mu$.
\end{corollary}

\begin{proof}
A conjugacy $\phi_\mu$ transports Koenigs charts, $\tilde\Rr=\phi\circ\Rr$ and
$\tilde S=\phi\circ S$.  So $\tilde T=T$ up to the normalising translations.  Take a
base point inside the domain of $\phi$ and apply
Proposition~\ref{prop:gen-base}.  The parameter $p$ depends only on the
multipliers.
\end{proof}

\begin{remark}[Numerical illustration]\label{rem:gen-numerics}
The script \texttt{generality\_check.py} implements the construction of this
section for any family given by $f$ and its Taylor coefficients at the fixed
points.  The table lists the values at $\lambda=0.97$, from runs over $\lambda=0.5,\dots,0.97$; the fourth column tends to $\operatorname{Re}\kappa^{(1)}$ as $p\to0$.
\begin{center}\small
\begin{tabular}{lllll}
\toprule
family & $\rho$ & $|B_1|$ & $\operatorname{Re}\log(\tau_1/B_1\eee^{2\pi\ii a})/p$ & $|\tau_1|/|B_1|-1$\\
\midrule
$b^w$ & $1/3$ & $0.0890584364122$ & $-0.0101990$ & $-9.4\cdot10^{-6}$\\
$w^2+c$ & $1$ & $0.0597942361690$ & $-0.0150481$ & $-1.4\cdot10^{-5}$\\
$v-\sin v+c$ & $1$ & $0.0723786416245$ & $-0.0142133$ & $-1.3\cdot10^{-5}$\\
$w+(w^2-\epsilon)\eee^w$ & $0$ & $2744.09379131$ & $+1.08542$ & $+1.0\cdot10^{-3}$\\
\bottomrule
\end{tabular}
\end{center}
In each family $\tau_n$, $t_n/t_1^n$ and $\log(\tau_1/B_1\eee^{2\pi\ii a})/p$ converge at
the rate $O(p)$.  The branch $\operatorname{Im}t_0=h/2$ holds to $10^{-31}$ or
better, and the sewing iteration of Theorem~\ref{thm:W2} gives
$\hat c_n=\tau_n(1+O(\Lambda))$.  For $w^2+c$ with the base points $w_0=0$, $0.2$ and
$0.35$, the real part in the fourth column agrees to all ten digits shown at each $\lambda$, as
Proposition~\ref{prop:gen-base} requires.  The imaginary part ranges from
$-0.07$ to $116$.  The germs of $w^2+c$ and $v-\sin v+c$ have the same
$\rho$ but different $|B_1|$ and $\operatorname{Re}\kappa^{(1)}$.
\end{remark}

\begin{remark}[Numerical evidence for the analogue of Theorem~\ref{mthm:F} for $w^2+c$]\label{rem:gen-quad}
For $f_c(w)=w^2+c$ the parameter $c=1/4$ is the cusp of the main cardioid, which
plays the role of the Shell--Thron region.  For real $c>1/4$ the two fixed points
are complex conjugate and repelling, and the real axis escapes.  The
two-fixed-point construction of Kneser type converges at the sampled
parameters with the normalisation $F(0)=f_c^3(0)$, corresponding to the intended
critical-point marking $F(-3)=0$ (script \texttt{quad\_kneser.py}).
The separate proof document \texttt{docs/proofs/quadratic-continuation.tex}
establishes its quadratic F and G using marked tilted domains, an explicit
polynomial multiplier deformation with open-overlap comparison, and an
independently replayed 4,013-box Arb exterior cover. Its conclusion concerns
regular height germs with this normalisation; it does not imply convergence
of the present numerical iteration. The following earlier observations
provide numerical evidence only.
\begin{itemize}
\item Along the arc $c=\tfrac14+0.1\eee^{\ii\varphi}$ the sampled values are
  consistent with continuation across the cardioid boundary.  Interpolation
  through points on both sides of the
  neutral band predicts held-out points to $2\cdot10^{-8}$, as accurately as
  interpolation within one side.
\item On vertical ladders $c=c_r+\ii y$, extrapolation towards $y=0$ supports
  agreement with the sewn coefficient.  At $\lambda=0.1$ the extrapolated value gives
  $(\hat c_1/\tau_1-1)/\Lambda=0.2272-0.1120\,\ii$, against $0.2279-0.1115\,\ii$
  from the first-order formula at the end of Section~\ref{sec:W2}.  So the
  data support agreement with $\Kk^W$ including its $O(\Lambda)$ correction;
  they do not prove existence or identification of the boundary limit.  The ratio
  $c_2/c_1^2$ agrees with $t_2/t_1^2$ to $4\cdot10^{-8}$.
\end{itemize}
For the analogue of Theorem~\ref{mthm:G} the numerics are only partial.  The
construction converges at sampled points in a wedge to the lower right of the
cusp.  In the multiplier coordinate $c=\mu/2-\mu^2/4$, at
$\arg\mu=0.25\pi$ and $0.4\pi$, interpolation across the neutral band errs by
$\le10^{-7}$ against a variation of order $10^{-1}$.  Further left the iteration
does not converge, and above it converges only after the $\theta$ lines are
raised in some tested cases.  These observations are consistent with local
holomorphic continuation, but do not establish it. Moving the Taylor centre
or fixing the constant Fourier mode of $\theta_{\mathrm{up}}$ may enlarge
the regular numerical region; these changes remain untested. They cannot
alone give direct coverage of the whole upper half-plane while retaining
both Schr\"oder charts: at $c=1/4+i/2$, $\lambda_{\mathrm{up}}=i$, the
fifth inverse-Schr\"oder coefficient requires $0=(3-5i)/2$.
The separate quadratic proof document supplies the exact obstruction and
a conditional local convergence theorem. A full numerical method must
also replace the neutral-end representation.
A pointwise algorithm that does not record paths does not by itself establish
single-valued analytic continuation; parameter holomorphy and agreement of
local branches must also be established.
\end{remark}
